% !TeX root = ../papers/nc-p3.tex
\section{非整な正則特殊化に必要な四次または八次の中心零因子}
\label{sec:finite-short-trace-pole}
\begin{theorem}\label{thm:finite-short-trace-pole}
十五次元成分の閉包で四辺が有限な標準正則窓を取る。合成整域でなければ
経路長4または8の非零斉次中心元 $c$ が存在し、任意に指定した頂点隅で $c_i y=0$ となる非零斉次元 $y$ がある。
相手 $y$ の経路長にも一様上界があり、例えば
\[
                       1860+8\,4^{1860}7^{255}
\]
で足りる。従って一つの隅でこの二因子の有限条件を排除した候補では、84次検査が正則性そのものと同値である。
この大きい上界は存在証明であり、全係数で消去を実施したという主張ではない。
\end{theorem}
\begin{proof}[特殊化先の相対trace環における八次正規化]
一般窓から指定点へ離散付値環 $O$ で特殊化し、$t$ をパラメータとする。一般も特殊も正則なので
各経路格子は自由であり、折り畳みを $\mathcal B$ と書く。一般の全行列次数16の被約traceを $\tau$ とする。
一般中心内で経路の全traceが生成する $O$ 代数を $C$、$T=\mathcal B C$ とする。
定理\ref{thm:noetherian-pi-window-cutoff}の一般行列からの特殊化により $C$ は有限型で $T$ は $C$ 上有限である。
ここでは頂点冪等元と四箇所の矢を一般分裂体の行列として代入し、係数環を $O$ に延長すれば同じ有限生成系が使える。
$C,T$ は各次数で有限自由な $O$ 加群で、一般繊維では通常の全中心と折り畳みに一致する。

$C/tC$ の非頂点閉指標で得るtraceは、特殊窓の全頂点次元4の半単純表現のtraceである。
実際 $\operatorname{Spec} C$ の一般点からこの指標へ付値環で近づけ、一般表現を分裂させた後に
$T$ 安定な自由格子を選べる。頂点射影はその階数4を保存する。残余体の代数閉拡大を行ってよい。
一般の絶対単純表現がこの半単純型へ特殊化するため、特殊繊維が両向き第一Extで直交する
二種の組成因子に分かれることはない。組成列を分けてから非対角項を厚みごとに消す証明は
定理\ref{thm:pi-vertex-extension-criterion}と同じで、関係の係数が変形しても前段まで対角なら
その係数変化は最初の非対角コサイクルに現れない。ここでの相対trace環が特殊窓自身の中心に
埋め込まれるとはまだ仮定していない。

非頂点の単純因子は薄いか、全頂点が次元2以上である。一頂点で次元4に達するものは全頂点が4である。
実際各隅は一般隅からの恒等式を満たす。四行四列行列の非恒等な中心多項式（Formanekの中心多項式）を、
その位置で全行列環を張る正長経路の線型結合に代入し、非零にできる。正次数の斉次中心成分が非零に作用し、
命題\ref{prop:all-window-corner-center}で全中心に延びる。従って頂点単純との全Extが零、四次差分のEuler式から
頂点次元は一定である。この場合は定理\ref{thm:index-four-octic-normalization}の行列論証が使え、
四文字または八文字のtraceが非零である。

残りでも四文字traceと八文字traceがともに全て消えることはない。頂点単純は正長のtraceに寄与しない。
非薄因子が重複込みで二つなら双方が $2\delta$ であり、定理\ref{thm:index-four-quartic-basepoints}の
復元と逆符号対の八次非消滅を使える。一つだけで薄い因子が二つなら $2\delta+\delta+\delta$ で、
四次の隣接切断階数4が薄い二つでは打ち消せない。最小頂点次元が3の非薄因子一つなら全て3で、
薄い因子がなければ三行三列の二段積次元5以上、あれば同じ定理の一文字階数の議論で四次traceは消えない。
最小頂点次元が2で薄い因子が高々一つの場合は、最小頂点を始点としてその非薄因子の四段積の線型像を
$L\subset M_2$ とする。各非零ベクトルが最後の辺で少なくとも二次元を生成するので $L$ は全非零ベクトル上全射、
従って $\dim L\ge3$。薄い因子がなければ四次・八次の消滅は $L$ が $\mathfrak{sl}_2$ 内で全等方であることを意味し矛盾する。
薄い因子が一つなら四次消滅によりその評価は $X\in L$ に対して $-\operatorname{tr}X$ であり、
八次消滅は $\operatorname{tr}(XY)+\operatorname{tr}(X)\operatorname{tr}(Y)=0$ となる。
これは $M_2$ 上非退化な形式で、全等方次元は高々2なので再び矛盾する。
全て薄い場合、各台点で倍率の一次・二次の重み付き和がともに零となる必要がある。
異なる高々四台点の四文字評価は独立で、八文字も単元となる四文字切断を掛ければ独立である。
各台点で少なくとも三つの異なる倍率型が必要なので一台点だけであり、異倍率同士と頂点との第一Ext消滅に反する。
全てが頂点単純なら正次数traceが全て零で、指定指標が非頂点であることに反する。

以上より $C/tC$ の経路長8は射影化で基点を持たない（四次traceの平方も含める）。
一般四元を取り、$C$ の八次の元に持ち上げると $C$ はそれらが生成する $O$ 部分環上有限である。
特殊繊維で四元を割れば有限次元であり、その有限基底の持上げが全次数で張ることは
次数別Nakayamaと中心次数の帰納から従う。
\end{proof}

\begin{proof}[短い極の吸収と相手の一様次数]
四文字および八文字の全traceが $\mathcal B$ の全頂点格子で整なら、$C$ の八次片も全て整である。
そこは八文字traceと四文字traceの積で張られるからである。このとき前段の四元の作るNoether環上で
$T$ は有限、部分加群 $\mathcal B$ も有限である。従って特殊窓は中心有限、正則性から合成整域である。
そこで非整域なら、長さ4または8の一語 $b$ で極を持つ $z=\tau(b)$ がある。
$t^u z$ を全頂点で整にする最小整数 $u>0$ を取り、剰余を $c$ とすれば非零中心元である。
任意の指定隅でも $c_i\ne0$ である。全中心から隅への制限の単射を使う。
その隅の一般traceとしても $z$ を読める。$b$ をその位置まで巡回した語の被約traceが同じだからである。

一般隅では、四行四列に対する交代多項式 $F$（定理\ref{thm:noetherian-pi-window-cutoff}）の
経路値 $f\ne0$ を長さ1860以下で取れる。中心体上、長さ4のループの累積積空間は安定すれば全隅になり、
単位から始めて高々15個の積で十六次元全体を張る。16個の交代引数と15個の間の引数をこれらの語から選ぶ。
交代最上外冪のtrace挿入式により $f z^m$ は全ての $m\ge0$ で指定隅の経路格子に入る。
$f$ の格子内の付値を $j$ とすれば $y_0=\overline{t^{-j}f}\ne0$ であり、
$(t^uz)^m t^{-j}f=t^{um-j}f z^m$ を還元して $c_i^{j+1} y_0=0$ を得る。
消える直前の元を $y$ と取れば長さは $1860+8j$ 以下である。

最後に付値環の選択に依存しない上界を与える。有限窓テンソルの十五次元成分に対応する
$\mathbb P^{255}$ の既約成分 $X$ を使う。ここでは基底変更前の射影テンソルを取る。
二つの隣接切断の七次小行列で定まる軌跡の成分なので $\deg X\le7^{255}$ である。
例えば $X$ の余次元個の七次式の一般線型結合でその一般点で孤立した成分にしてBézoutを適用する。
先に一般点で上の $f$ を選び、特殊点を通る一般超平面切断曲線の一つの枝を使う。
その次数は $\deg X$ 以下で、一般点は必要な稠密開集合にあり、$f$ の非消滅も保つ。
実際避けるべき真閉部分との一般線型切断は指定点の近くで高々零次元である。
曲線の正規化の分枝の離散付値環を取り、テンソルの一つの非零座標で局所自明化する。
長さ $l\le1860$ の固定経路空間は $4^l$ 行の表示を持ち、関係列は四テンソルの隣接縮約を語の間に入れたものなので
テンソル座標に一次である。特殊点と一般点で関係の階数 $r$ が同じなので特殊点で単元の $r$ 次小行列を選べる。
一定列 $f$ を加えた行列の非零な $(r+1)$ 次小行列を取れば、その零点次数は $r\deg X$ 以下である。
実際 $f$ が剰余自由格子内で $t^j$ 倍ならこの全ての小行列が $t^j$ で割れる。列 $f$ が定数であるから
斉次化した小行列の次数は $r$ であり、投影自明化の分母は指定点で単元である。
よって $j\le4^{1860}7^{255}$。主張した有限範囲で零因子がなければ、正則点は必ず整域なので84次判定に戻る。
\end{proof}
